% ============================================================
%  CHAPITRE : ÉTUDE D'UNE FONCTION
%  Mazava Maths — Terminale S (Bacc Série S, Madagascar)
%  Aligné sur le modèle type (Isométries) — 2026-08-27.
%  Contenu recentré (demande utilisateur) :
%    - plan d'étude d'une fonction
%    - diagramme d'étude des branches à l'infini
%    - convexité / concavité / point d'inflexion (théorème + propriété
%      + application)
%    - position relative d'une courbe et d'une droite
%    - identification de f, f', f'' sur une même figure
%  Parti pris : tous les exemples et exercices utilisent ln et exp.
%  Environnements : definition, propriete, theoreme, methode,
%                   attention, astucebac, exemple, exercice,
%                   solution, remarque
% ============================================================

\chapter{Étude d'une fonction}

\begin{citationchapitre}
	« Toute la connaissance d'une courbe tient dans ses pentes et sa
	courbure. »\\[0.4em]
	— d'après Christiaan Huygens
\end{citationchapitre}

\begin{objectifs}
	\begin{itemize}[label=\textbullet]
		\item Conduire l'\textbf{étude complète} d'une fonction selon un plan
		systématique et tracer sa courbe.
		\item Déterminer la nature d'une \textbf{branche infinie} (asymptote
		horizontale, verticale, oblique ; direction asymptotique ; branche
		parabolique) à l'aide d'un \textbf{diagramme}.
		\item Étudier la \textbf{convexité}, la \textbf{concavité}, déterminer les
		\textbf{points d'inflexion} ; en déduire des inégalités.
		\item Étudier la \textbf{position relative} d'une courbe et d'une droite
		(asymptote, tangente, sécante) ; résoudre graphiquement.
		\item Reconnaître les courbes de $f$, $f'$ et $f''$ sur une \textbf{même
		figure}.
	\end{itemize}
\end{objectifs}

\begin{remarque}
	La continuité, le TVI, le calcul des dérivées et la convexité (théorème
	$f''\geq 0\iff f$ convexe) ont été vus aux chapitres \og Limites et
	continuité \fg\ et \og Dérivabilité \fg. On les réinvestit ici pour
	\emph{étudier} et \emph{tracer} des courbes.
\end{remarque}


% ================================================================
\section{Plan d'étude d'une fonction}
% ================================================================

\begin{methode}
	\textbf{Plan d'étude complète d'une fonction $f$.}
	\begin{enumerate}
		\item \textbf{Ensemble de définition} $D_f$ : exclure les valeurs
		annulant un dénominateur ou rendant $<0$ l'argument d'un $\sqrt{\;}$ ou
		$\leq 0$ celui d'un $\ln$.
		\item \textbf{Parité / périodicité} (si $D_f$ est symétrique par rapport
		à $0$) : réduire l'intervalle d'étude.
		\item \textbf{Limites aux bornes} de $D_f$ : en déduire les
		\textbf{branches infinies} (section suivante).
		\item \textbf{Dérivée} $f'(x)$ : calcul, puis \textbf{signe} (souvent via
		une fonction auxiliaire).
		\item \textbf{Tableau de variations} : signe de $f'$, sens de variation,
		valeurs et limites aux bornes, extremums.
		\item \textbf{Convexité} : signe de $f''$, points d'inflexion.
		\item \textbf{Points particuliers} : intersections avec les axes,
		tangentes remarquables, position par rapport aux asymptotes.
		\item \textbf{Tracé} : asymptotes, points et tangentes clés, allure.
	\end{enumerate}
\end{methode}

\begin{propriete}{Signe de $f'$ et variations}{prop_signe_fprime}
	Soit $f$ dérivable sur un intervalle $I$.
	\begin{center}
	\renewcommand{\arraystretch}{1.5}
	\begin{tabular}{>{\centering\arraybackslash}p{4.6cm}
	                >{\centering\arraybackslash}p{4.2cm}}
	\hline
	\textbf{Sur $I$} & \textbf{Conséquence} \\
	\hline
	$f'>0$          & $f$ strictement croissante \\
	$f'<0$          & $f$ strictement décroissante \\
	$f'=0$          & $f$ constante \\
	$f'$ change de signe en $a$ & extremum local en $a$ \\
	\hline
	\end{tabular}
	\end{center}
\end{propriete}

\exerciceresolu
\textbf{Étude complète de $f(x)=(x-1)\e^{x}$}

\solution

\textbf{1. Domaine.} $D_f=\R$. Pas de parité.

\textbf{2. Limites.}
\[
\lim_{x\to-\infty}(x-1)\e^{x}=\lim_{x\to-\infty}\bigl(x\e^{x}-\e^{x}\bigr)=0,
\qquad
\lim_{x\to+\infty}(x-1)\e^{x}=+\infty.
\]
La droite $y=0$ est \textbf{asymptote horizontale} en $-\infty$.

\textbf{3. Dérivée.} $f'(x)=\e^{x}+(x-1)\e^{x}=x\,\e^{x}$, du signe de $x$.

\textbf{4. Variations.}
\begin{center}
\renewcommand{\arraystretch}{1.4}
\begin{tabular}{c|ccccc}
$x$      & $-\infty$ & & $0$ & & $+\infty$ \\
\hline
$f'(x)$  &  & $-$ & $0$ & $+$ & \\
\hline
$f(x)$   & $0$ & $\searrow$ & $-1$ & $\nearrow$ & $+\infty$ \\
\end{tabular}
\end{center}
Minimum global $f(0)=-1$.

\textbf{5. Convexité.} $f''(x)=\e^{x}+x\e^{x}=(x+1)\e^{x}$, du signe de
$x+1$. $f$ est concave sur $]-\infty\,;-1]$, convexe sur $[-1\,;+\infty[$ :
\textbf{point d'inflexion} en $\bigl(-1\,;-2\e^{-1}\bigr)$.

\textbf{6. Points particuliers.} $f(x)=0\iff x=1$ ; $f(0)=-1$.
En $+\infty$, $\dfrac{f(x)}{x}=\dfrac{(x-1)\e^{x}}{x}\to+\infty$ :
\textbf{branche parabolique} de direction $(Oy)$.

\begin{center}
\begin{tikzpicture}[>=Stealth, scale=0.95]
	\draw[->] (-4.2,0) -- (2.2,0) node[right,font=\small] {$x$};
	\draw[->] (0,-1.6) -- (0,3.0) node[above,font=\small] {$y$};
	\draw[thick,domain=-4.1:1.35,samples=140] plot (\x,{(\x-1)*exp(\x)});
	\filldraw (0,-1) circle (1.3pt) node[right,font=\scriptsize] {$(0\,;-1)$};
	\filldraw (1,0) circle (1.3pt) node[above right,font=\scriptsize] {$1$};
	\filldraw (-1,{-2*exp(-1)}) circle (1.3pt)
	  node[left,font=\scriptsize] {infl.};
	\node[font=\scriptsize] at (1.15,2.3) {$\mathcal C_f$};
\end{tikzpicture}
\end{center}


% ================================================================
\section{Diagramme d'étude des branches infinies}
% ================================================================

\begin{definition}
	Une \textbf{branche infinie} de $\mathcal C_f$ est une portion de courbe
	qui s'éloigne indéfiniment de l'origine, lorsque $x\to a$ (avec
	$|f(x)|\to+\infty$) ou lorsque $x\to\pm\infty$.
\end{definition}

\begin{definition}
	Soit $\mathcal C_f$ la courbe de $f$ et $\Delta$ une droite.
	\begin{itemize}[label=\textbullet]
		\item $\Delta:\;y=\ell$ est \textbf{asymptote horizontale} en $\pm\infty$
		si $\displaystyle\lim_{x\to\pm\infty}f(x)=\ell$.
		\item $\Delta:\;x=a$ est \textbf{asymptote verticale} si
		$\displaystyle\lim_{x\to a}f(x)=\pm\infty$ (éventuellement d'un seul côté).
		\item $\Delta:\;y=ax+b$ ($a\neq 0$) est \textbf{asymptote oblique} en
		$\pm\infty$ si $\displaystyle\lim_{x\to\pm\infty}\bigl[f(x)-(ax+b)\bigr]=0$.
		\item Si $\dfrac{f(x)}{x}\to a\in\R^*$ mais $f(x)-ax\to\pm\infty$, la
		courbe admet la \textbf{direction asymptotique} de la droite $y=ax$
		(branche parabolique dans cette direction).
		\item Si $\dfrac{f(x)}{x}\to 0$ (resp.\ $\pm\infty$) avec
		$f(x)\to\pm\infty$, la courbe admet une \textbf{branche parabolique} de
		direction $(Ox)$ (resp.\ $(Oy)$).
	\end{itemize}
\end{definition}

\begin{center}
\begin{tikzpicture}[>=Stealth, scale=0.62, transform shape,
	every node/.style={font=\small},
	box/.style={draw, rounded corners, align=center, inner sep=4pt,
	            fill=gray!8, minimum height=8mm},
	res/.style={draw, rounded corners, align=center, inner sep=4pt,
	            fill=black!78, text=white, minimum height=8mm}]

	\node[box] (l1) at (0,0) {$\displaystyle\lim_{x\to\pm\infty}f(x)$ ?};
	\node[res] (ah) at (7.4,0) {Asymptote\\ horizontale $y=\ell$};
	\draw[->] (l1) -- (ah) node[midway,above,font=\scriptsize] {$=\ell\in\R$};

	\node[box] (l2) at (0,-2.6) {$\displaystyle\lim_{x\to\pm\infty}\dfrac{f(x)}{x}$ ?};
	\draw[->] (l1) -- (l2) node[midway,left,font=\scriptsize] {$=\pm\infty$};

	\node[res] (box) at (7.4,-1.8) {Branche parabolique\\ direction $(Ox)$};
	\draw[->] (l2) -- (box) node[midway,above,font=\scriptsize] {$=0$};
	\node[res] (boy) at (7.4,-3.4) {Branche parabolique\\ direction $(Oy)$};
	\draw[->] (l2) -- (boy) node[midway,below,font=\scriptsize] {$=\pm\infty$};

	\node[box] (l3) at (0,-5.6) {$\displaystyle\lim_{x\to\pm\infty}\bigl[f(x)-ax\bigr]$ ?};
	\draw[->] (l2) -- (l3) node[midway,left,font=\scriptsize] {$=a\in\R^*$};

	\node[res] (obl) at (7.4,-5.0) {Asymptote oblique\\ $y=ax+b$};
	\draw[->] (l3) -- (obl) node[midway,above,font=\scriptsize] {$=b\in\R$};
	\node[res] (dir) at (7.4,-6.6) {Direction asymptotique\\ $y=ax$ (br.\ parabolique)};
	\draw[->] (l3) -- (dir) node[midway,below,font=\scriptsize] {$=\pm\infty$};
\end{tikzpicture}
\end{center}

\begin{attention}
	Pour $x\to a$ (fini) : si $\displaystyle\lim_{x\to a}f(x)=\pm\infty$, alors
	$x=a$ est asymptote verticale ; si la limite est finie, il n'y a pas de
	branche infinie en $a$ (au plus un \og trou \fg).
\end{attention}

\begin{methode}
	Recherche d'une \textbf{asymptote oblique} $y=ax+b$ en $+\infty$ :
	\begin{enumerate}
		\item $a=\displaystyle\lim_{x\to+\infty}\dfrac{f(x)}{x}$ ;
		\item si $a$ est un réel non nul, $b=\displaystyle\lim_{x\to+\infty}
		\bigl[f(x)-ax\bigr]$ ;
		\item si $b$ est fini, $y=ax+b$ est asymptote oblique.
	\end{enumerate}
	Astuce : si $f(x)=ax+b+\varphi(x)$ avec $\varphi(x)\to 0$, l'asymptote se
	lit directement et $\varphi(x)$ donne la \textbf{position}.
\end{methode}

\begin{exemple}
	\begin{itemize}[label=\textbullet]
		\item $f(x)=\ln x$ : $\dfrac{\ln x}{x}\to 0$ et $\ln x\to+\infty$
		$\Rightarrow$ branche parabolique de direction $(Ox)$.
		\item $f(x)=\e^{x}$ : $\dfrac{\e^{x}}{x}\to+\infty$ $\Rightarrow$ branche
		parabolique de direction $(Oy)$ (et asymptote $y=0$ en $-\infty$).
		\item $f(x)=x+\e^{-x}$ : $\dfrac{f(x)}{x}\to 1$ et $f(x)-x=\e^{-x}\to 0$
		$\Rightarrow$ asymptote oblique $y=x$.
		\item $f(x)=x+\ln x$ : $\dfrac{f(x)}{x}\to 1$ mais $f(x)-x=\ln x\to+\infty$
		$\Rightarrow$ direction asymptotique $y=x$, sans asymptote.
	\end{itemize}
\end{exemple}


% ================================================================
\section{Convexité, concavité et points d'inflexion}
% ================================================================

\begin{definition}
	Soit $f$ dérivable sur un intervalle $I$.
	\begin{itemize}[label=\textbullet]
		\item $f$ est \textbf{convexe} sur $I$ si $\mathcal C_f$ est entièrement
		\textbf{au-dessus de chacune de ses tangentes} (de façon équivalente :
		en dessous de chacune de ses cordes).
		\item $f$ est \textbf{concave} sur $I$ si $\mathcal C_f$ est en dessous de
		chacune de ses tangentes.
	\end{itemize}
\end{definition}

\begin{theoreme}[Caractérisation de la convexité]
	Soit $f$ deux fois dérivable sur $I$.
	\[
	f \text{ convexe sur } I
	\iff f'' \geq 0 \text{ sur } I
	\iff f' \text{ croissante sur } I.
	\]
	($f$ concave $\iff f''\leq 0 \iff f'$ décroissante.)
\end{theoreme}

\begin{propriete}{Point d'inflexion}{prop_inflexion}
	$A\bigl(a\,;f(a)\bigr)$ est un \textbf{point d'inflexion} de $\mathcal C_f$
	lorsque $\mathcal C_f$ \textbf{traverse} sa tangente en $A$. Si $f$ est deux
	fois dérivable, cela équivaut à :
	\[
	f''(a)=0 \quad\text{et}\quad f'' \text{ change de signe en } a.
	\]
	($f''(a)=0$ seul ne suffit pas : ex.\ $f(x)=x^{4}$ en $0$, avec le repère
	adapté ; ici on prendra plutôt des exemples $\ln$/$\exp$.)
\end{propriete}

\begin{center}
\begin{tikzpicture}[>=Stealth, scale=0.9]
	% convexe
	\begin{scope}
		\draw[->] (-0.3,0) -- (3.2,0) node[right,font=\scriptsize] {$x$};
		\draw[->] (0,-0.3) -- (0,2.6);
		\draw[thick,domain=-0.1:1.55,samples=60] plot (\x,{exp(\x)-0.2});
		\draw[blue,dashed,domain=0.1:1.9] plot (\x,{exp(0.9)*(\x-0.9)+exp(0.9)-0.2});
		\node[font=\scriptsize] at (1.6,-0.35) {convexe : $f''\geq 0$};
		\node[font=\scriptsize,blue] at (2.0,0.55) {tangente};
	\end{scope}
	% concave
	\begin{scope}[xshift=5.2cm]
		\draw[->] (-0.3,0) -- (3.2,0) node[right,font=\scriptsize] {$x$};
		\draw[->] (0,-0.3) -- (0,2.6);
		\draw[thick,domain=0.12:3.0,samples=60] plot (\x,{1.4*ln(\x)+1.1});
		\draw[red,dashed,domain=0.3:2.6] plot (\x,{(1.4/1.4)*(\x-1.4)+1.4*ln(1.4)+1.1});
		\node[font=\scriptsize] at (1.6,-0.35) {concave : $f''\leq 0$};
		\node[font=\scriptsize,red] at (2.6,2.15) {tangente};
	\end{scope}
\end{tikzpicture}
\end{center}

\begin{propriete}{Inégalités de convexité}{prop_ineg_convexite}
	Puisque $\mathcal C_f$ est du même côté de chacune de ses tangentes :
	\begin{itemize}[label=\textbullet]
		\item $\exp$ est convexe : $\mathcal C$ au-dessus de la tangente en $0$
		($y=1+x$), donc $\;\e^{x}\geq 1+x\;$ pour tout $x$ ;
		\item $\ln$ est concave : $\mathcal C$ en dessous de la tangente en $1$
		($y=x-1$), donc $\;\ln x\leq x-1\;$ pour tout $x>0$.
	\end{itemize}
\end{propriete}

\begin{methode}
	Pour \textbf{étudier la convexité} de $f$ :
	\begin{enumerate}
		\item calculer $f''(x)$ ;
		\item étudier le signe de $f''$ (tableau) ;
		\item conclure : intervalles de convexité / concavité, points d'inflexion
		(changement de signe de $f''$).
	\end{enumerate}
\end{methode}

\exerciceresolu
\textbf{Convexité de $f(x)=x-\ln(x^{2}+1)$}

Étudier la convexité de $f$ sur $\R$ et déterminer les points d'inflexion.

\solution

\[
f'(x)=1-\frac{2x}{x^{2}+1}
=\frac{x^{2}-2x+1}{x^{2}+1}
=\frac{(x-1)^{2}}{x^{2}+1}.
\]
En dérivant ce quotient :
\[
\begin{aligned}
f''(x)
&=\frac{2(x-1)(x^{2}+1)-(x-1)^{2}\cdot 2x}{(x^{2}+1)^{2}} \\[0.3em]
&=\frac{2(x-1)\bigl[(x^{2}+1)-x(x-1)\bigr]}{(x^{2}+1)^{2}}
 =\frac{2(x-1)(x+1)}{(x^{2}+1)^{2}}.
\end{aligned}
\]

Le signe de $f''$ est celui de $(x-1)(x+1)$ :
\begin{center}
\renewcommand{\arraystretch}{1.4}
\begin{tabular}{c|ccccccc}
$x$       & $-\infty$ & & $-1$ & & $1$ & & $+\infty$ \\
\hline
$f''(x)$  & & $+$ & $0$ & $-$ & $0$ & $+$ & \\
\hline
$\mathcal C_f$ & & convexe & & concave & & convexe & \\
\end{tabular}
\end{center}
$f''$ change de signe en $-1$ et en $1$ : $\mathcal C_f$ admet \textbf{deux
points d'inflexion}, d'abscisses $-1$ et $1$.


% ================================================================
\section{Position relative d'une courbe et d'une droite}
% ================================================================

\begin{methode}
	Pour étudier la position de $\mathcal C_f$ par rapport à la droite
	$\Delta:\;y=mx+p$ :
	\begin{enumerate}
		\item poser $\varphi(x)=f(x)-(mx+p)$ ;
		\item étudier le \textbf{signe} de $\varphi(x)$ ;
		\item conclure :
		\begin{itemize}[label=\textbullet]
			\item $\varphi(x)>0$ : $\mathcal C_f$ est \textbf{au-dessus} de $\Delta$ ;
			\item $\varphi(x)<0$ : $\mathcal C_f$ est \textbf{en dessous} de $\Delta$ ;
			\item $\varphi(x)=0$ : \textbf{point commun} (intersection, ou contact
			si $\varphi$ ne change pas de signe).
		\end{itemize}
	\end{enumerate}
\end{methode}

\begin{attention}
	Trois situations fréquentes :
	\begin{itemize}[label=\textbullet]
		\item $\Delta$ \textbf{asymptote} : $\varphi(x)\to 0$ en $\pm\infty$, et
		le signe de $\varphi$ donne la position \emph{au voisinage} de l'infini ;
		\item $\Delta$ \textbf{tangente} en $a$ : $\varphi(a)=\varphi'(a)=0$
		($\varphi$ admet $a$ pour racine \og double \fg) ;
		\item $\Delta$ \textbf{sécante} : $\varphi$ change de signe en chaque
		point d'intersection.
	\end{itemize}
\end{attention}

\begin{exemple}
	$f(x)=\e^{x}$ et sa tangente en $0$, $\Delta:\;y=x+1$.
	$\varphi(x)=\e^{x}-x-1$ ; $\varphi'(x)=\e^{x}-1$ s'annule en $0$ en
	changeant de signe : $\varphi$ admet un minimum $\varphi(0)=0$. Donc
	$\varphi(x)\geq 0$ : $\mathcal C_f$ est \textbf{au-dessus} de $\Delta$, avec
	contact en $(0\,;1)$.
\end{exemple}

\begin{exemple}
	$f(x)=x+\e^{-x}$ et son asymptote $\Delta:\;y=x$.
	$\varphi(x)=\e^{-x}>0$ : $\mathcal C_f$ est \textbf{toujours au-dessus} de
	$\Delta$.
\end{exemple}

\begin{astucebac}
	\textbf{Résolution graphique.} Résoudre $f(x)\leq mx+p$ revient à lire les
	$x$ pour lesquels $\mathcal C_f$ est \emph{en dessous} de $\Delta$, i.e.
	$\varphi(x)\leq 0$. De même, $f(x)=k$ se lit à l'intersection de
	$\mathcal C_f$ avec l'horizontale $y=k$, et $f(x)=g(x)$ à l'intersection de
	$\mathcal C_f$ et $\mathcal C_g$.
\end{astucebac}

\begin{exemple}
	Résoudre $\e^{x}\leq x+2$. On pose $\varphi(x)=\e^{x}-x-2$.
	$\varphi'(x)=\e^{x}-1$ : $\varphi$ décroît sur $]-\infty\,;0]$, croît
	ensuite, $\varphi(0)=-1<0$. De plus $\varphi(-2)=\e^{-2}>0$ et
	$\varphi(1)=\e-3<0$… en fait $\varphi(x)\to+\infty$ aux deux bornes.
	Par le TVI, $\varphi$ s'annule une fois sur $]-\infty\,;0[$ (en un réel
	$\alpha<0$) et une fois sur $]0\,;+\infty[$ (en $\beta>0$). L'ensemble des
	solutions est $[\alpha\,;\beta]$.
\end{exemple}


% ================================================================
\section{Reconnaître \texorpdfstring{$f$}{f}, \texorpdfstring{$f'$}{f'} et \texorpdfstring{$f''$}{f''} sur une même figure}
% ================================================================

\begin{methode}
	Pour identifier trois courbes $\mathcal C_1,\mathcal C_2,\mathcal C_3$
	représentant $f$, $f'$, $f''$ (dans le désordre) :
	\begin{enumerate}
		\item là où \textbf{$f$ a un extremum}, $f'$ \textbf{s'annule} en
		changeant de signe ;
		\item là où \textbf{$f$ croît}, $f'$ est \textbf{positive} ; là où $f$
		décroît, $f'<0$ ;
		\item là où \textbf{$f$ a un point d'inflexion}, $f''$ \textbf{s'annule}
		en changeant de signe (et $f'$ y a un extremum) ;
		\item là où \textbf{$f$ est convexe}, $f''>0$ et $f'$ est croissante ;
		\item comparer les comportements aux bornes (une exponentielle domine sa
		dérivée d'un facteur constant ; dériver \og aplatit \fg\ souvent les
		bosses).
	\end{enumerate}
	On identifie d'abord une paire \og fonction / dérivée \fg, puis on
	recommence avec la troisième.
\end{methode}

\exerciceresolu
\textbf{Identification sur la figure}

La figure représente, dans le désordre, $f$, $f'$ et $f''$ avec
$f(x)=\e^{-x^{2}}$.

\begin{center}
\begin{tikzpicture}[>=Stealth, scale=1.15]
	\draw[->] (-2.4,0) -- (2.4,0) node[right,font=\scriptsize] {$x$};
	\draw[->] (0,-2.3) -- (0,1.4) node[above,font=\scriptsize] {$y$};
	\foreach \x in {-2,-1,1,2} \draw (\x,-0.05)--(\x,0.05);
	% f
	\draw[very thick,domain=-2.2:2.2,samples=120] plot (\x,{exp(-\x*\x)});
	\node[font=\scriptsize] at (0.42,1.05) {$\mathcal C_1$};
	% f'
	\draw[thick,dash pattern=on 3pt off 2pt,domain=-2.2:2.2,samples=120]
	  plot (\x,{-2*\x*exp(-\x*\x)});
	\node[font=\scriptsize] at (-1.15,1.05) {$\mathcal C_2$};
	% f''
	\draw[thick,dotted,domain=-2.2:2.2,samples=140]
	  plot (\x,{(4*\x*\x-2)*exp(-\x*\x)});
	\node[font=\scriptsize] at (0.3,-1.95) {$\mathcal C_3$};
\end{tikzpicture}
\end{center}

\begin{enumerate}
	\item Identifier $\mathcal C_1$, $\mathcal C_2$, $\mathcal C_3$.
	\item Justifier à l'aide des points où les courbes coupent l'axe $(Ox)$.
\end{enumerate}

\solution

\textbf{1.} $\mathcal C_1$ est une \og cloche \fg\ toujours positive, de
maximum $1$ en $0$ : c'est $f$. \\
$\mathcal C_2$ est impaire, s'annule en $0$ (là où $f$ a son maximum),
négative pour $x>0$ (là où $f$ décroît) : c'est $f'$. \\
$\mathcal C_3$ est paire, vaut $-2$ en $0$, positive \og aux ailes \fg :
c'est $f''$.

\textbf{2.} $\mathcal C_2$ coupe $(Ox)$ en $x=0$, abscisse du maximum de
$\mathcal C_1$ : cohérent avec $f'=0$ aux extremums de $f$.
$\mathcal C_3$ coupe $(Ox)$ en $x=\pm\dfrac{1}{\sqrt 2}$, abscisses des
points d'inflexion de $\mathcal C_1$ et des extremums de $\mathcal C_2$ :
cohérent avec $f''=(f')'$.


% ================================================================
\section{Méthodes et exercices résolus}
% ================================================================

\begin{methode}
	\textbf{Nature d'une branche infinie en $+\infty$ — mode d'emploi rapide.}
	\begin{center}
	\renewcommand{\arraystretch}{1.4}
	\begin{tabular}{>{\raggedright\arraybackslash}p{4.4cm}
	                >{\raggedright\arraybackslash}p{3.8cm}}
	\hline
	\textbf{On observe} & \textbf{Conclusion} \\
	\hline
	$f(x)\to \ell$          & asymptote $y=\ell$ \\
	$f(x)\to\pm\infty$, $\dfrac{f(x)}{x}\to 0$
	                        & br.\ parabolique dir.\ $(Ox)$ \\
	$\dfrac{f(x)}{x}\to\pm\infty$
	                        & br.\ parabolique dir.\ $(Oy)$ \\
	$\dfrac{f(x)}{x}\to a$, $f(x)-ax\to b$
	                        & asymptote $y=ax+b$ \\
	$\dfrac{f(x)}{x}\to a$, $f(x)-ax\to\pm\infty$
	                        & direction asymptotique $y=ax$ \\
	\hline
	\end{tabular}
	\end{center}
\end{methode}

\exerciceresolu
\textbf{Branches infinies}

Préciser la nature de la branche infinie en $+\infty$ de :
\[
\text{a) } f(x)=\frac{2\e^{x}+1}{\e^{x}+1}
\qquad
\text{b) } g(x)=x+2-\ln x
\qquad
\text{c) } h(x)=\frac{x^{2}}{\e^{x}}
\]

\solution

\textbf{a)} $f(x)=\dfrac{2+\e^{-x}}{1+\e^{-x}}\to 2$ : asymptote horizontale
$y=2$.

\textbf{b)} $\dfrac{g(x)}{x}=1+\dfrac{2}{x}-\dfrac{\ln x}{x}\to 1$, mais
$g(x)-x=2-\ln x\to-\infty$ : direction asymptotique $y=x$, pas d'asymptote.

\textbf{c)} $h(x)=\dfrac{x^{2}}{\e^{x}}\to 0$ : asymptote horizontale $y=0$.

\bigskip
\exerciceresolu
\textbf{Position par rapport à une asymptote}

Soit $f(x)=x-1+\dfrac{2}{\e^{x}}$. Montrer que $\Delta:\;y=x-1$ est
asymptote oblique en $+\infty$ et étudier la position de $\mathcal C_f$ par
rapport à $\Delta$.

\solution

$f(x)-(x-1)=2\,\e^{-x}\xrightarrow[x\to+\infty]{}0$ : $\Delta:\;y=x-1$ est
asymptote oblique en $+\infty$.

$\varphi(x)=2\,\e^{-x}>0$ pour tout $x$ : $\mathcal C_f$ est \textbf{au-dessus}
de $\Delta$ sur $\R$.

\bigskip
\exerciceresolu
\textbf{Une inégalité par convexité}

Montrer que pour tout $x>0$, $\;\ln x\leq \dfrac{x}{\e}$.

\solution

Soit $g(x)=\ln x$, concave sur $]0\,;+\infty[$ (car $g''(x)=-\tfrac{1}{x^{2}}<0$).
$\mathcal C_g$ est donc en dessous de chacune de ses tangentes. La tangente
au point d'abscisse $\e$ a pour pente $g'(\e)=\dfrac{1}{\e}$ et passe par
$(\e\,;1)$ : son équation est $y=\dfrac{1}{\e}(x-\e)+1=\dfrac{x}{\e}$.
Donc $\ln x\leq \dfrac{x}{\e}$ pour tout $x>0$ (égalité en $x=\e$).


% ================================================================
\section{Exercices d'entraînement}
% ================================================================

% --------------------------------------------------
\subsection{Échauffement : lire une figure}
% --------------------------------------------------

\exercice{Lire une courbe et sa dérivée \hfill \facile}
On donne ci-dessous les courbes $\mathcal C_1$ et $\mathcal C_2$ : l'une
représente une fonction $f$, l'autre sa dérivée $f'$.

\begin{center}
\begin{tikzpicture}[>=Stealth, scale=1.0]
	\draw[->] (-2.6,0) -- (2.6,0) node[right,font=\scriptsize] {$x$};
	\draw[->] (0,-1.4) -- (0,2.0) node[above,font=\scriptsize] {$y$};
	\foreach \x in {-2,-1,1,2} \draw (\x,-0.05)--(\x,0.05);
	\draw[very thick,domain=-2.3:1.4,samples=120] plot (\x,{exp(\x)-1});
	\node[font=\scriptsize] at (1.25,1.9) {$\mathcal C_1$};
	\draw[thick,dash pattern=on 3pt off 2pt,domain=-2.3:1.1,samples=120]
	  plot (\x,{exp(\x)});
	\node[font=\scriptsize] at (0.95,1.35) {$\mathcal C_2$};
\end{tikzpicture}
\end{center}

\begin{enumerate}
	\item Laquelle des deux courbes représente $f$ ? Laquelle représente $f'$ ?
	Justifier par le signe.
	\item $\mathcal C_2$ coupe-t-elle l'axe $(Ox)$ ? Qu'en déduit-on pour les
	variations de $f$ ?
	\item Que valent $\displaystyle\lim_{x\to-\infty}f(x)$ et
	$\displaystyle\lim_{x\to-\infty}f'(x)$ ?
	\item $f$ est-elle convexe ? Justifier à partir de $\mathcal C_2$.
\end{enumerate}

\exercice{Diagramme des branches infinies \hfill \facile}
Pour chaque fonction, suivre le diagramme et conclure sur la branche
infinie en $+\infty$.
\begin{multicols}{2}
\begin{enumerate}
	\item $f(x)=\dfrac{\e^{x}}{\e^{x}+2}$
	\item $f(x)=\ln(x+1)$
	\item $f(x)=x+\dfrac{\ln x}{x}$
	\item $f(x)=\e^{x}-x$
	\item $f(x)=\dfrac{x+\ln x}{x}$
	\item $f(x)=x-\e^{-x}$
\end{enumerate}
\end{multicols}

\exercice{Vrai ou faux \hfill \facile}
Répondre par \emph{vrai} ou \emph{faux} en justifiant brièvement.
\begin{enumerate}
	\item Si $\displaystyle\lim_{x\to+\infty}\dfrac{f(x)}{x}=2$, alors
	$\mathcal C_f$ admet une asymptote oblique.
	\item Si $f''$ s'annule en $a$, alors $\mathcal C_f$ a un point d'inflexion
	en $a$.
	\item Si $f$ est convexe, alors $\mathcal C_f$ est au-dessus de chacune de
	ses tangentes.
	\item Si $\mathcal C_f$ et $\Delta:\;y=x$ vérifient $f(x)-x>0$ pour tout $x$,
	alors $\Delta$ n'est pas asymptote.
	\item La courbe de $f'$ coupe l'axe des abscisses aux abscisses des
	extremums de $f$.
	\item $\ln x\leq x-1$ pour tout $x>0$.
	\item Une fonction dont la dérivée est croissante est convexe.
	\item $\e^{x}$ et $\ln x$ ont une asymptote oblique en $+\infty$.
\end{enumerate}

% --------------------------------------------------
\subsection{Branches infinies et convexité}
% --------------------------------------------------

\exercice{Nature des branches infinies \hfill \moyen}
Étudier \emph{toutes} les branches infinies (bornes de $D_f$ et $\pm\infty$).
\begin{multicols}{2}
\begin{enumerate}
	\item $f(x)=\dfrac{\ln x}{x}$
	\item $f(x)=x\,\e^{-x}$
	\item $f(x)=x+1+\e^{-x}$
	\item $f(x)=\dfrac{\e^{x}}{x}$
	\item $f(x)=x-\ln x$
	\item $f(x)=\dfrac{1}{\ln x}$
\end{enumerate}
\end{multicols}

\exercice{Convexité \hfill \moyen}
Étudier la convexité et déterminer les points d'inflexion.
\begin{multicols}{2}
\begin{enumerate}
	\item $f(x)=x\,\e^{x}$
	\item $f(x)=(x^{2}-2)\e^{x}$
	\item $f(x)=x-\ln(x^{2}+1)$
	\item $f(x)=\dfrac{\e^{x}}{\e^{x}+1}$
\end{enumerate}
\end{multicols}

\exercice{Inégalités par convexité \hfill \moyen}
\begin{enumerate}
	\item En utilisant la tangente à $\mathcal C_{\exp}$ en $0$, montrer que
	$\e^{x}\geq 1+x$ pour tout $x$.
	\item En utilisant la tangente à $\mathcal C_{\ln}$ en $1$, montrer que
	$\ln x\leq x-1$ pour tout $x>0$.
	\item En déduire que pour tout $x>0$, $\;1-\dfrac1x\leq \ln x\leq x-1$.
\end{enumerate}

\exercice{Position relative \hfill \moyen}
Soit $f(x)=\dfrac{\e^{x}-1}{\e^{x}+1}$ et $\Delta:\;y=\dfrac{x}{2}$.
\begin{enumerate}
	\item Montrer que $f$ est impaire.
	\item Étudier $\varphi(x)=f(x)-\dfrac{x}{2}$ : calculer $\varphi'(x)$ et
	$\varphi''(x)$, puis dresser le signe de $\varphi'$.
	\item En déduire la position de $\mathcal C_f$ par rapport à $\Delta$.
\end{enumerate}

% --------------------------------------------------
\subsection{Identifier \texorpdfstring{$f$}{f}, \texorpdfstring{$f'$}{f'}, \texorpdfstring{$f''$}{f''}}
% --------------------------------------------------

\exercice{Trois courbes \hfill \moyen}
La figure représente $f$, $f'$ et $f''$ dans le désordre, avec
$f(x)=(x-1)\e^{x}$.

\begin{center}
\begin{tikzpicture}[>=Stealth, scale=1.0]
	\draw[->] (-3.4,0) -- (1.7,0) node[right,font=\scriptsize] {$x$};
	\draw[->] (0,-1.7) -- (0,2.2) node[above,font=\scriptsize] {$y$};
	\foreach \x in {-3,-2,-1,1} \draw (\x,-0.05)--(\x,0.05);
	\draw[very thick,domain=-3.3:1.05,samples=140] plot (\x,{(\x-1)*exp(\x)});
	\node[font=\scriptsize] at (0.95,-1.4) {$\mathcal C_1$};
	\draw[thick,dash pattern=on 3pt off 2pt,domain=-3.3:0.9,samples=140]
	  plot (\x,{\x*exp(\x)});
	\node[font=\scriptsize] at (0.75,1.1) {$\mathcal C_2$};
	\draw[thick,dotted,domain=-3.3:0.75,samples=140]
	  plot (\x,{(\x+1)*exp(\x)});
	\node[font=\scriptsize] at (0.35,1.9) {$\mathcal C_3$};
\end{tikzpicture}
\end{center}

\begin{enumerate}
	\item Vérifier que $f'(x)=x\,\e^{x}$ et $f''(x)=(x+1)\e^{x}$.
	\item Associer chaque courbe à $f$, $f'$ ou $f''$.
	\item Retrouver l'abscisse du minimum de $\mathcal C_1$ et celle de son
	point d'inflexion à partir des zéros de $\mathcal C_2$ et $\mathcal C_3$.
\end{enumerate}

\exercice{Lecture croisée \hfill \moyen}
On sait que $g'$ est représentée par la courbe ci-dessous.
\begin{center}
\begin{tikzpicture}[>=Stealth, scale=1.0]
	\draw[->] (-0.4,0) -- (4.3,0) node[right,font=\scriptsize] {$x$};
	\draw[->] (0,-1.3) -- (0,1.6) node[above,font=\scriptsize] {$y$};
	\foreach \x in {1,2,3,4} \draw (\x,-0.05)--(\x,0.05) node[below,font=\tiny]{$\x$};
	\draw[very thick,domain=0.35:4.1,samples=120] plot (\x,{ln(\x)-0.7});
	\node[font=\scriptsize] at (3.6,1.0) {$\mathcal C_{g'}$};
\end{tikzpicture}
\end{center}
\begin{enumerate}
	\item Résoudre graphiquement $g'(x)=0$ ; en déduire l'abscisse de
	l'extremum de $g$ et sa nature.
	\item Donner le sens de variation de $g$.
	\item $\mathcal C_g$ est-elle convexe ou concave ? Justifier.
\end{enumerate}

% --------------------------------------------------
\subsection{Sujets types Bac}
% --------------------------------------------------

\exercice{Sujet type Bac — étude complète de $(x+1)\e^{-x}$ \hfill \difficile}
Soit $f(x)=(x+1)\,\e^{-x}$, de courbe $\mathcal C$.

\textbf{Partie A — Étude}
\begin{enumerate}
	\item Déterminer $\displaystyle\lim_{x\to-\infty}f(x)$ et
	$\displaystyle\lim_{x\to+\infty}f(x)$ ; préciser les branches infinies.
	\item Montrer que $f'(x)=-x\,\e^{-x}$ et dresser le tableau de
	variations ; préciser le maximum.
	\item Montrer que $f''(x)=(x-1)\,\e^{-x}$, étudier la convexité et
	déterminer le point d'inflexion $I$ de $\mathcal C$.
	\item Déterminer l'équation de la tangente $T$ à $\mathcal C$ au point
	d'abscisse $0$.
\end{enumerate}

\textbf{Partie B — Position et tracé}
\begin{enumerate}
	\setcounter{enumi}{4}
	\item Étudier la position de $\mathcal C$ par rapport à $T$ au voisinage
	de $0$.
	\item Résoudre $f(x)=0$ et déterminer les coordonnées des points
	d'intersection de $\mathcal C$ avec les axes.
	\item Tracer $T$, l'asymptote $y=0$, le point $I$ et la courbe
	$\mathcal C$.
\end{enumerate}

\exercice{Sujet type Bac — la fonction $\dfrac{\ln x}{x}$ \hfill \difficile}
Soit $f(x)=\dfrac{\ln x}{x}$ sur $]0\,;+\infty[$, de courbe $\mathcal C$.

\textbf{Partie A — Branches infinies et variations}
\begin{enumerate}
	\item Étudier les branches infinies de $\mathcal C$ (en $0^{+}$ et en
	$+\infty$).
	\item Calculer $f'(x)$ et dresser le tableau de variations ; préciser le
	maximum.
	\item Calculer $f''(x)$ et montrer que $\mathcal C$ a un unique point
	d'inflexion, d'abscisse $\e^{3/2}$.
\end{enumerate}

\textbf{Partie B — Résolution graphique}
\begin{enumerate}
	\setcounter{enumi}{3}
	\item Déterminer l'image de $]0\,;+\infty[$ par $f$.
	\item Discuter, selon le réel $m$, le nombre de solutions de l'équation
	$\ln x=m\,x$.
	\item Donner l'équation de la tangente à $\mathcal C$ au point d'abscisse
	$1$ et étudier la position locale de $\mathcal C$ par rapport à elle.
\end{enumerate}

\exercice{Sujet type Bac — asymptote oblique et position \hfill \difficile}
Soit $f(x)=x-2+\dfrac{\ln x}{x}$ sur $]0\,;+\infty[$, de courbe $\mathcal C$.

\textbf{Partie A}
\begin{enumerate}
	\item Déterminer $\displaystyle\lim_{x\to 0^{+}}f(x)$ ; en déduire une
	asymptote.
	\item Montrer que la droite $\Delta:\;y=x-2$ est asymptote à $\mathcal C$
	en $+\infty$.
	\item Étudier la position de $\mathcal C$ par rapport à $\Delta$.
\end{enumerate}

\textbf{Partie B}
\begin{enumerate}
	\setcounter{enumi}{3}
	\item Calculer $f'(x)$. On admet que $f'(x)$ a le signe de
	$g(x)=x^{2}+1-\ln x$. Étudier $g$ et en déduire le signe de $f'$.
	\item Dresser le tableau de variations de $f$.
	\item Montrer que l'équation $f(x)=0$ admet une unique solution
	$\alpha$, et que $1<\alpha<2$.
\end{enumerate}

\exercice{Sujet type Bac — identification $f$, $f'$, $f''$ \hfill \difficile}
Soit $f(x)=\ln\bigl(1+x^{2}\bigr)$, de courbe $\mathcal C$.

\textbf{Partie A — Étude}
\begin{enumerate}
	\item Étudier la parité de $f$, puis $\displaystyle\lim_{x\to\pm\infty}f(x)$
	et la nature de la branche infinie.
	\item Montrer que $f'(x)=\dfrac{2x}{1+x^{2}}$ ; dresser le tableau de
	variations.
	\item Montrer que $f''(x)=\dfrac{2\bigl(1-x^{2}\bigr)}{\bigl(1+x^{2}\bigr)^{2}}$
	et que $\mathcal C$ a deux points d'inflexion, d'abscisses $\pm 1$.
\end{enumerate}

\textbf{Partie B — Trois courbes}
Sur la figure ci-dessous sont tracées $f$, $f'$ et $f''$, notées
$\mathcal C_1$, $\mathcal C_2$, $\mathcal C_3$ dans le désordre.

\begin{center}
\begin{tikzpicture}[>=Stealth, scale=1.0]
	\draw[->] (-3.3,0) -- (3.3,0) node[right,font=\scriptsize] {$x$};
	\draw[->] (0,-1.0) -- (0,2.6) node[above,font=\scriptsize] {$y$};
	\foreach \x in {-3,-2,-1,1,2,3} \draw (\x,-0.05)--(\x,0.05);
	\draw[very thick,domain=-3.1:3.1,samples=140]
	  plot (\x,{ln(1+\x*\x)});
	\node[font=\scriptsize] at (2.7,2.15) {$\mathcal C_1$};
	\draw[thick,dash pattern=on 3pt off 2pt,domain=-3.1:3.1,samples=140]
	  plot (\x,{2*\x/(1+\x*\x)});
	\node[font=\scriptsize] at (1.15,1.2) {$\mathcal C_2$};
	\draw[thick,dotted,domain=-3.1:3.1,samples=160]
	  plot (\x,{2*(1-\x*\x)/((1+\x*\x)*(1+\x*\x))});
	\node[font=\scriptsize] at (0.55,-0.75) {$\mathcal C_3$};
\end{tikzpicture}
\end{center}

\begin{enumerate}
	\setcounter{enumi}{3}
	\item Identifier $\mathcal C_1$ (courbe toujours positive, sans maximum).
	\item Identifier $\mathcal C_2$ et $\mathcal C_3$ en utilisant : $\mathcal C_2$
	s'annule là où $\mathcal C_1$ a son minimum, et $\mathcal C_3$ s'annule là où
	$\mathcal C_2$ a ses extremums.
	\item Justifier que la courbe de $f'$ est impaire et celle de $f''$ paire.
\end{enumerate}

\exercice{Sujet type Bac — famille $f_m(x)=x+m\,\e^{-x}$ \hfill \difficile}
Pour $m\in\R$, on pose $f_m(x)=x+m\,\e^{-x}$, de courbe $\mathcal C_m$.

\textbf{Partie A — La famille}
\begin{enumerate}
	\item Montrer que la droite $\Delta:\;y=x$ est asymptote à toutes les
	$\mathcal C_m$ en $+\infty$, et étudier la position de $\mathcal C_m$ par
	rapport à $\Delta$ selon le signe de $m$.
	\item Calculer $f_m'(x)$ et $f_m''(x)$. Étudier, selon $m$, l'existence
	d'un extremum et sa nature.
	\item Pour $m>0$, montrer que le minimum de $f_m$ a pour coordonnées
	$\bigl(\ln m\,;\,\ln m+1\bigr)$ ; en déduire que ce minimum décrit la
	droite $y=x+1$ quand $m$ décrit $]0\,;+\infty[$.
\end{enumerate}

\textbf{Partie B — Le cas $m=1$}
\begin{enumerate}
	\setcounter{enumi}{3}
	\item Étudier $f_1(x)=x+\e^{-x}$ : variations, convexité.
	\item Montrer que $\mathcal C_1$ est au-dessus de $\Delta$ et que $f_1$
	admet un minimum égal à $1$ en $x=0$.
	\item Tracer $\Delta$ et $\mathcal C_1$.
\end{enumerate}
